\scanpage{311}
\mspage{157}

The set $\unsure{\widetilde U|U}^{!}$ of points $Q\in U$ such that $G_Q^+\ne(1)$ is stable under the action of $\mathcal G$, and is identified (with this action) with \unsure{the set} of \inserted{maximal} liftings modulo $\pi$ of (cyclic) subgroups $\ne1$ of $G^{+}_{\deleted{\unsure{Q}}}$. When one knows, for a partial lifting $r$ of a $G_Q$ into $\mathcal E$ what its centralizer in $\mathcal E$ is i.e. the corresponding lifting of $\mathcal G_Q$, \unsure{well|then} in the $=$ way for the conjugates of $r$ by any \unsure{element} $g$ (not only of $\pi$, but $=$ of $\mathcal E$), by simple conjugation.

So the cases to be determined correspond in practice to the orbits of $\mathcal G$ in $U^{!}$.

One may be interested in describing $\mathcal G$ or $\Gamma$ as a subgroup of $\mathrm{Aut\,ext}_{\mathrm{lac}}(\pi)=\Gamma'$ \inserted{giving rise to the extension $\mathcal E'$ of $\Gamma'$ by $\pi$ (thus $\mathcal E\subset\mathcal E'$).} But for every lifting $r$ of a $G_Q$, let us consider $\mathrm{Centr}_{\mathcal E'}(r)$, \deleted{\illegible{1}} . \unsure{we shall have|we still have} $\mathrm{Centr}_{\mathcal E'}(r)\cap\pi=(1)$ i.e. \inserted{we find a section over} the image of this centralizer in $\Gamma'$, say $\Gamma'(Q)$. This image depends only on $Q$ i.e. on the $\pi$-conjugacy class of $r$ or of $r(\unsure{G_Q})$, and is replaced by a $G$-conjugate when $Q$ is replaced by a $G$-conjugate. This

\mspage{310}

\scanpage{312}

said, the intersection \inserted{$\Gamma'^{\natural}$} of the $\Gamma'(Q)$, for $Q\in U^{!}$, is a subgroup of $\Gamma'$ which contains the intersection \deleted{\unsure{of}} $\mathcal G^{\natural}$ of the $\mathcal G_Q^{\deleted{\unsure{\natural}}}$, and the centralizer of $G$ in $\Gamma'^{\natural}$ contains in the $=$ way the intersection \deleted{of} $\Gamma^{\natural}$ of the $\Gamma_Q$, which is an \inserted{invariant} subgroup of finite index in $\Gamma$. And one can then propose to oneself whether it is possible to characterize at least \unsure{the|a} finite subgroup $\Gamma^{\natural}$ of $\Gamma$ as $\mathrm{Centr}_{\Gamma'^{\natural}}(G)$, and \unsure{to recover from it} perhaps $\Gamma$ as the normalizer of $\Gamma^{\natural}$ in \deleted{the centralizer} $\mathrm{Centr}_{G}(\Gamma)$.

I am more particularly interested in the profinite variant of this, in the case where $U=\mathbb P^1_{\mathbb Q}\setminus\{0,1,\infty\}$, $G=\mathfrak S_3$, $\Gamma={}$ Galois group of $\overline{\mathbb Q}/\mathbb Q$, \deleted{while} \inserted{$\unsure{P|\rho}=\exp(2i\pi/6)$}.
\margin{algebraic closure in $\mathbb C$}

I have not checked that $\Gamma\longrightarrow\mathrm{Autext}_{\mathrm{lac}}(\widehat\pi)$ is injective, this does not prevent doing calculations in $\mathrm{Aut}_{\mathrm{lac}}(\widehat\pi)$, be they heuristic for the moment.

\mspage{311}

\scanpage{313}
\mspage{158}

I am running into notational troubles -- too many groups are \unsure{\uline{denoted}} by the letter $\Gamma$ (possibly with primes, subscripts, superscripts \ldots). There are three types of groups that enter into my reflections:

a) the Teichmüller groups and its variants, which play the role \unsure{of the|of} ``universal'' groups operating (\inserted{possibly} modulo isotopy) on surfaces, or on outer groups \deleted{\unsure{with loops}} with loops. These groups tend to be infinite

b) Groups (most often finite) operating on topological surfaces, or on geometric algebraic curves (\unsure{without}, in that case, moving the base field).

\margin{Cases b) and c) sometimes mix (in \unsure{a} \unsure{$G$} \unsure{groups} $\mathcal G$ extension of a Galois group $\Gamma$ by a finite group \unsure{$G$}) \unsure{in the case of} \unsure{$G$} Algebraic \unsure{Geometry}.}

c) Profinite Galois groups (hence infinite), chiefly of fields of finite type over $\mathbb Q$, operating ``arithmetically'' on surfaces and their geometric $\widehat{\pi}_1$.

It is because of the deep analogies between cases b) and c), and their close relations with case a), that I had been led to adopt common notations, but

\mspage{312}

\scanpage{314}

which in the long run end up leading to collisions. There is thus reason to review the notations. I will reserve the letter $G$ \inserted{and variants} for geometric actions (case b)) of groups, most often finite, the letter $\Gamma$ \inserted{and variants} for Galois groups, the letter $\mathcal G$ for mixed groups.

As for the ``universal'' groups of Teichmüller type, like those denoted $\Gamma_{g,\nu}$ previously, I will instead denote them \inserted{$\mathcal T$,} $\mathcal T_{g,\nu}$ (initial of ``Teichmüller'', whereas $\Gamma$, \inserted{$G$} \unsure{are} the initial of Galois).

The Galois group \inserted{over $\mathbb Q$} of the algebraic closure $\overline{\mathbb Q}$ of $\mathbb Q$ in $\mathbb C$ \deleted{\unsure{of $\overline{\mathbb Q}$ over}} \unsure{deserves} a special letter, I will denote it $\boldsymbol{\Gamma}$. The quotient \deleted{\unsure{by} $\widehat{\Gamma}_{g,\nu}$ of} $\operatorname{Norm}_{\widehat{\widehat{\mathcal T}}_{g,\nu}}(\widehat{\mathcal T}_{g,\nu})/\widehat{\mathcal T}_{g,\nu}$, which is more akin to a Galois group than to a Teichmüller group, will be denoted $\boldsymbol{\Gamma}_{g,\nu}$ (bold letter!). In the

\mspage{313}

\scanpage{315}
\mspage{159}

case $(g,\nu)=(0,3)$ which occupies me more particularly, $\boldsymbol{\Gamma}_{0,3}$ is identified with the \deleted{\unsure{normalizer}} \inserted{centralizer} of $G=\mathfrak S_3$ \inserted{$=\mathcal T_{0,3}^{+}$} in $\widehat{\widehat{\mathcal T}}_{0,3}$ -- it is contained in $\widehat{\widehat{\mathcal T}}_{03}^{!}$, \inserted{and perhaps equal.} We have a canonical homomorphism
\margin{$\widehat{\widehat{\mathcal T}}$ is here the $\frac{1}{2}$ direct product of $\mathfrak S_3=\mathcal T^+$ by $\mathcal T^{!}$}
\[
\boldsymbol{\Gamma}\longrightarrow\boldsymbol{\Gamma}_{03}
\qquad (\text{more generally }\boldsymbol{\Gamma}\longrightarrow\boldsymbol{\Gamma}_{g,\nu})
\]
for which I do not know for the moment whether it is injective, and still less whether it is surjective. The preceding reflections suggest conditions on the image, which are especially interesting if one admits the relations
\[
\widehat{\pi}^{\rho}=\widehat{\pi}^{\sigma_0}=(1)
\]

We \inserted{have} \unsure{denoted\deleted{will}} by $\mathcal E_{g\nu}$ the canonical extension of $\mathcal T_{g,\nu}$ by $\pi_{g,\nu}$ (which for $(g,\nu)$ anabelian is identified with $\mathcal T'_{g,\nu+1}$).

The outer operation of a group $G$, $\Gamma$ or $\mathcal G$ also defines an extension by $\pi$ (or by $\widehat{\pi}$), which we have also \unsure{denoted|denoted} by the letter $\mathcal E$ (initial of extension) -- there is again a collision of notations.

\mspage{314}

\scanpage{316}

I will take the letter $\mathfrak{S}$ (which is a counterpart to $\mathcal{T}$) for these extensions in the universal Teichmüller-style cases, (writing $\mathfrak{S}_{g,\nu}$ instead of the $\Gamma'_{g,\nu+1}$, since the viewpoint is different\ldots), and keeping the letter $\mathcal{E}$ in \unsure{the|the} \unsure{preceding|preceding} cases. So $\mathcal{E}$ tends to be a subextension of an $\mathfrak{S}$.

Admitting that $\boldsymbol{\Gamma}\longrightarrow\boldsymbol{\Gamma}_{0,3}$ is injective, we would thus have

\diagram{\begin{tikzcd}[ampersand replacement=\&,column sep=small,row sep=large]
1 \arrow[r,no head] \& \widehat{\pi}_{0,3} \arrow[r] \arrow[d,equal] \& \widehat{\widehat{\mathfrak{S}}}_{0,3} \arrow[r] \arrow[d] \& \widehat{\widehat{\mathcal{T}}}_{0,3} \arrow[r] \arrow[d] \& 1 \\
1 \arrow[r,no head] \& \widehat{\pi}_{0,3} \arrow[r] \arrow[d,equal] \& \widehat{\widehat{\mathfrak{S}}}'_{0,3} \arrow[r] \arrow[d] \& \underbrace{\unsure{\mathcal{T}^{+}_{0,3}}}_{\mathfrak{S}_{3}}\;\deleted{\unsure{\boldsymbol{\Gamma}}}\!\times\boldsymbol{\Gamma}_{0,3} \arrow[r] \arrow[d] \& 1 \\
1 \arrow[r] \& \widehat{\pi}_{03} \arrow[r] \arrow[d] \& \mathcal{E} \arrow[r] \arrow[d] \& \mathfrak{S}_{3}\times\boldsymbol{\Gamma} \arrow[r,no head] \arrow[d] \& 1 \\
1 \arrow[r] \& \pi_{0,3} \arrow[r] \& \mathfrak{S}_{03} \arrow[r,no head] \& \underbrace{\mathfrak{S}_{3}\times\unsure{\mathbb{Z}/2}}_{\mathcal{T}_{0,3}} \arrow[r] \& 1
\end{tikzcd}}

\margin{N.B. We denote by $\widehat{\widehat{\mathfrak{S}}}'_{g,\nu}$ the normalizer of $\widehat{\mathfrak{S}}^{+}_{g,\nu}$ in $\widehat{\widehat{\mathfrak{S}}}_{g,\nu}$, an extension of $\boldsymbol{\Gamma}_{g,\nu}$ by $\widehat{\mathfrak{S}}^{+}_{g,\nu}$}

\unsure{If|Let} $\gamma\in\boldsymbol{\Gamma}_{0,3}$, \deleted{\unsure{which is}}\inserted{for it to be} in the image of $\boldsymbol{\Gamma}$ \deleted{\unsure{have|is} the property}\inserted{it is necessary} \unsure{that it|that they} \unsure{admit|admit} a lifting $u$ which commutes with $\rho$, and a lifting which commutes with $\sigma_{0}$ \unsure{(which|Which}, \unsure{as soon|of the} as $\gamma\in\widehat{\widehat{\mathcal{T}}}_{0,3}$, implies

\mspage{315}

\scanpage{317}
\mspage{160}

already that $\gamma$ in $\widehat{\widehat{\mathcal{T}}}_{0,3}$ commutes with $\mathcal{T}^{+}_{0,3}=\mathfrak{S}_{3}$, \deleted{thus} i.e. that they is in $\boldsymbol{\Gamma}_{0,3}$). It could be that \unsure{every} element of $\boldsymbol{\Gamma}_{0,3}$ already has this property, hence that this condition does not \unsure{impose|make} any genuine restriction on the image of $\boldsymbol{\Gamma}$ in $\boldsymbol{\Gamma}_{0,3}$.

\mspage{316}

\scanpage{318}
\mspage{161}

\uline{26} Profinite Teichmüller groups.

\margin{Discretifications and prediscretifications.}

Let $\pi$ be a profinite group with loops \inserted{of type $g,\nu$}, $T$ the invertible $\widehat{\mathbb Z}$-module of its orientations.

We assume we are in the \unsure{anabelian} case, and we shall admit \unsure{that then}
\[
\operatorname{Centre}(\pi)=\{1\}
\tag{1}
\]
more generally that the centralizer in $\pi$ of every open subgroup of $\pi$ is reduced to $\{1\}$. We shall thus again have a canonical exact sequence of profinite groups
\[
1 \unsure{\longrightarrow|{-}} \pi \longrightarrow \operatorname{Aut}_{\mathrm{lac}}(\pi) \longrightarrow \operatorname{Autext}_{\mathrm{lac}}(\pi) \longrightarrow 1
\tag{\unsure{2}}
\]
We shall also put \deleted{\unsure{$\operatorname{Autext}$}}
\[
\operatorname{Autext}_{\mathrm{lac}}(\pi)=\widehat{\widehat{\mathcal T}}(\pi)
\tag{3}
\]
\unsure{we|and} will call it the \uline{extended Teichmüller group} of $\pi$. We shall also put $\widehat{\widehat{\mathfrak S}}(\pi)=\operatorname{Aut}_{\mathrm{lac}}(\pi)$, so that we can write (2) as
\[
1 \longrightarrow \pi \longrightarrow \widehat{\widehat{\mathfrak S}}(\pi) \longrightarrow \widehat{\widehat{\mathcal T}}(\pi) \longrightarrow 1
\tag{2'}
\]

Let us call a ``basis'' of $\pi$ a \unsure{set} of \unsure{elements} of $\pi$ $(x_i,y_i)_{1\leq i\leq g}$, $(\ell_j)_{1\leq j\leq\nu}$, \inserted{the $\ell_j$ \unsure{generating$^{\mathrm t}$} the different groups with loops} satisfying
\[
\deleted{\unsure{\prod_i}\,\unsure{[x_i,y_i]}\,\unsure{\cdots[x_g,y_g]\ell_1}}
\quad
\inserted{\ell_\nu\ell_{\nu-1}\cdots\ell_1[x_g,y_g]\cdots[x_1,y_1]=1}
\tag{\unsure{4}}
\]
and such that this is a generating relation.

\mspage{317}

\scanpage{319}

\margin{or again, one defines $\Pi_{g,\nu}$ as the discrete group \unsure{with|with the} generators the $x_i,y_i,\ell_j$ and with \unsure{defining} relation (4)}

\unsure{If} one chooses in $\Pi_{g,\nu}$ a (discrete) basis (corresponding \unsure{definition|Definition}), \unsure{whence} a basis of $\widehat{\Pi}_{g,\nu}$, one will have an obvious bijection
\[
\operatorname{Bases}(\pi) \xleftarrow{\sim} \operatorname{Isom}_{\mathrm{lac}}(\widehat{\Pi}_{g,\nu},\pi)
\tag{5}
\]

The set of bases of $\pi$ is a homogeneous set \unsure{in|under} \deleted{\illegible{1}} $\widehat{\widehat{\mathfrak S}}(\pi)=\operatorname{Aut}_{\mathrm{lac}}(\pi)$. \unsure{Let|If} the \deleted{\unsure{$\Pi_0$ d}} basis corresponds to an isomorphism $u:\widehat{\Pi}_{g,\nu}\longrightarrow\pi$, the group \inserted{$\Pi_0$} generated by the $x_i,y_i,\ell_j$ is none other than $\deleted{\unsure{\Pi_0}}\inserted{u}(\Pi_{g,\nu})$. \deleted{One of the} \inserted{the subgroups} of $\pi$ which can be obtained in this way are called the \uline{discretifications} of $\pi$. These form a homogeneous set under $\widehat{\widehat{\mathfrak S}}=\operatorname{Aut}_{\mathrm{lac}}(\pi)$, \inserted{canonically} isomorphic to the quotient of the right $\widehat{\widehat{\mathfrak S}}_{g,\nu}$-\deleted{\unsure{torsor}}\inserted{set} $\operatorname{Isom}_{\mathrm{lac}}(\widehat{\Pi}_{g,\nu},\pi)$ by $\mathfrak S_{g,\nu}$:
\[
\operatorname{Discrt}(\pi) \xleftarrow{\sim} \operatorname{Isom}_{\mathrm{lac}}(\widehat{\Pi}_{g,\nu},\pi)/\mathfrak S_{g,\nu}
\tag{6}
\]
The \unsure{automorphisms} \inserted{(profinite with loops)} of $\pi$ fixing \inserted{a discretification} $\Pi_0$ are identified with the \unsure{automorphisms} of the discrete group \inserted{with loops} $\Pi_0$.
\[
\operatorname{Aut}_{\mathrm{lac}}(\pi,\Pi_0)\simeq\operatorname{Aut}_{\mathrm{lac}}(\Pi_0)
\tag{7}
\]

\mspage{318}

\scanpage{320}
\mspage{162}

The bijection (5) puts on the set of Bases of $\pi$ a bitorsor structure under $\widehat{\widehat{\mathfrak S}}(\pi),\widehat{\widehat{\mathfrak S}}_{g,\nu}$, and \unsure{the|the} \unsure{corresponding|corresponding} \unsure{topologies|topology} makes it a profinite set. The set of discretifications of $\pi$, which is a quotient of the preceding set, inherits a quotient topology, which is none other than the quotient topology of the right-hand side of (6). \deleted{Except in the \unsure{exceptional} cases \unsure{$(g,\nu)=(0,3)$} \unsure{where} $\mathfrak S_{g,\nu}$} \inserted{This space is} not Hausdorff \inserted{(N.B. $\mathfrak S_{g,\nu}$ is \unsure{always} infinite)}, for the group $\mathfrak S_{g,\nu}\subset\widehat{\widehat{\mathfrak S}}_{g,\nu}$ is not closed. We denote by $\unsure{\widetilde{\operatorname{Discrt}}'|\operatorname{Discrt}'}(\pi)$ the associated Hausdorff topological space, identified with $\operatorname{Isom}_{\mathrm{lac}}(\widehat{\Pi}_{g,\nu},\pi)/\unsure{\overline{\mathfrak S}}_{g,\nu}$, where $\overline{\mathfrak S}_{g,\nu}$ denotes the closure of $\mathfrak S_{g,\nu}$ in $\widehat{\widehat{\mathfrak S}}_{g,\nu}$. We moreover have an obvious surjective homomorphism $\widehat{\mathfrak S}_{g,\nu}\longrightarrow\unsure{\widehat{\widehat{\mathfrak S}}|\overline{\mathfrak S}}_{g,\nu}$ whose image is $\overline{\mathfrak S}_{g,\nu}$, \unsure{we shall admit} that it is injective, and shall identify $\overline{\mathfrak S}_{g,\nu}$ with $\widehat{\mathfrak S}_{g,\nu}$. Thus
\[
\operatorname{Discrt}'(\pi)\simeq\operatorname{Isom}_{\mathrm{lac}}(\widehat{\Pi}_{g,\nu},\pi)/\widehat{\mathfrak S}_{g,\nu}
\tag{\unsure{8}}
\]

\mspage{319}